\section{The source language and visibility discipline}
\label{sec:source-language}

\subsection{Expression languages}

The core protocol syntax is parameterized by an expression language
\(L\).  An expression language supplies a type universe \(\Ty\), a
family of semantic values \(\Val(\tau)\), expression and distribution
syntax, evaluation functions, finite dependency sets, and soundness
laws saying that evaluation depends only on the declared dependencies.
The event-graph compiler treats these dependencies as declared read
footprints.  An instance may over-approximate, but it may not omit a
semantic dependency.

The surface language currently specializes to the simple expression
language used by the examples.  That specialization is deliberate:
nullable yield sugar needs options, default values, and a commitment
payload predicate that are concrete to that expression language.  The
maintained core remains generic in \(L\).

\subsection{Visibility-aware contexts}

A visibility-aware binding has the form \((x,\pub\,\tau)\) or
\((x,\hid\,p\,\tau)\).  A public binding is visible to every player; a
hidden binding is visible only to its owner.  From a context \(\Gamma\)
we use three projections:
\[
  \erasepub{\Gamma}, \qquad
  \erasectx{\Gamma}, \qquad
  \view{p}.
\]
The first keeps only public bindings and erases visibility.  The second
keeps all bindings and erases visibility.  The third keeps public
bindings plus bindings hidden to player \(p\).

Core syntax is indexed by the visibility context, so ordinary variable
lookup is proof-bearing.  Freshness and well-formed contexts make
lookup proof-irrelevant: two proofs naming the same variable in the
same checked context identify the same typed field during compilation.

\subsection{Core protocol syntax}

The maintained core has four constructors:
\[
\arraycolsep=1pt
\begin{array}{r@{\;}c@{\;}l}
\Ret\,\overline{p}
  & : & \mathsf{VegasCore}\,\Gamma, \\
\Sample\,x \sim D;\,k
  & : & \mathrm{DistExpr}(\erasepub{\Gamma}, b) \times
       \mathsf{VegasCore}((x,\pub\,b)::\Gamma)
       \to \mathsf{VegasCore}\,\Gamma, \\
\Commit_p\,x \mid R;\,k
  & : & \Expr((x,b)::\erasectx{\view{p}}, \bool) \times
       \mathsf{VegasCore}((x,\hid\,p\,b)::\Gamma)
       \to \mathsf{VegasCore}\,\Gamma, \\
\Reveal\,y := x;\,k
  & : & \mathrm{VHasVar}\,\Gamma\,x\,(\hid\,p\,b) \times
       \mathsf{VegasCore}((y,\pub\,b)::\Gamma)
       \to \mathsf{VegasCore}\,\Gamma.
\end{array}
\]
\(\Sample\) draws a public value.  \(\Commit\) asks player \(p\) to
choose a hidden value satisfying a guard over the proposed action and
the player's visible context.  \(\Reveal\) exposes a hidden value as a
fresh public alias.  \(\Ret\) terminates with payoff expressions over
the public context.

There is no core \(\Let\).  Deterministic administrative bindings live
in the surface layer and are erased by substitution before graph
compilation.  This keeps event graphs about protocol activity rather
than source convenience.  If a deterministic value must be materialized
as public protocol state, it should be expressed through a public sample
or an explicit future protocol constructor, not through administrative
\(\Let\).

\subsection{Surface syntax}
\label{sec:source:surface}

\(\mathsf{VegasLang}\) is the current surface language over the simple
expression language.  It mirrors the protocol core and adds:
\begin{itemize}\itemsep1pt
\item administrative \(\Let\)-bindings, erased by the lowering
  environment;
\item nullable \(\mathsf{yield}\), lowered to a hidden optional commit
  followed immediately by a public reveal of the optional value.
\end{itemize}

The lowering environment maps public surface variables to public core
expressions and hidden variables to hidden core variables.  This is why
surface lets do not allocate graph fields: later public references are
rewritten to the defining expression, while hidden commitments remain
stored values that can later be revealed.

\subsection{Static obligations}
\label{sec:source:wf}

The code distinguishes the obligations needed for graph construction
from the obligations needed for checked game semantics.

\begin{description}[itemsep=2pt,leftmargin=1.2em,style=nextline]
\item[\(\mathsf{GraphProgram}\).]
  Contains an initial visibility context, a core program, an initial
  environment, a well-formed-context proof, fresh-binder evidence, and
  normalized sample distributions.  This is the minimal input to
  \(\mathsf{compile}\).
\item[\(\mathsf{WFProgram}\).]
  Wraps a \(\mathsf{GraphProgram}\) and adds reveal completeness plus
  commit-guard legality.  Reveal completeness states that committed
  source values are eventually opened by the source program.  Legality
  states that each commit guard admits at least one action in every
  source-visible context.  These obligations are consumed by the
  game-facing progress layer, not by raw graph allocation.
\item[\(\mathsf{FiniteDomains}\).]
  Evidence that all operational value domains introduced by the program
  are finite.  This is needed to enumerate finite strategy carriers and
  finite machine states.  Terminal payoff integers need not themselves
  form a finite type.
\end{description}

\paragraph{Lean correspondence.}
Core syntax is \TT{Vegas.Core.Basic}.  Checked-program bundles are in
\TT{Vegas.Core.Program}.  Surface syntax and let-erasing lowering are
in \TT{Vegas.Lang.Basic}.  The concrete well-formedness predicates are
collected by \TT{Vegas.Core.WF}.
